\documentclass[12pt, a4paper]{article}
\usepackage[T2A]{fontenc}
\usepackage[utf8]{inputenc}
\usepackage[ukrainian]{babel}
\usepackage{amsmath, amssymb, amsthm}
\usepackage{geometry}
\geometry{top=2cm, bottom=2cm, left=2.5cm, right=1.5cm}
\usepackage{hyperref}

\title{\textbf{Основи теорії груп Лі, алгебр Лі та їх застосування}}
\author{}
\date{}

\begin{document}
	\maketitle
	
	\section{Що таке група та алгебра Лі?}
	
	\subsection{Поняття групи}
	У математиці \textbf{група} $G$ — це не просто набір об'єктів (наприклад, матриць або операторів), а множина, на якій задана операція множення, що підпорядковується чотирьом суворим правилам (аксіомам):
	\begin{enumerate}
		\item \textbf{Замкненість:} Добуток двох елементів групи також належить до цієї групи ($A, B \in G \implies A \cdot B \in G$). Наприклад, композиція двох поворотів — це знову поворот.
		\item \textbf{Асоціативність:} $(A \cdot B) \cdot C = A \cdot (B \cdot C)$.
		\item \textbf{Одиничний елемент:} Існує елемент $I$ (нейтральне перетворення, "нічого не робити"), такий що $A \cdot I = I \cdot A = A$.
		\item \textbf{Обернений елемент:} Для кожного перетворення $A$ існує зворотне $A^{-1}$, яке повертає систему в початковий стан: $A \cdot A^{-1} = I$.
	\end{enumerate}
	
	\textbf{Група Лі} — це група, елементи якої неперервно залежать від одного або декількох параметрів (наприклад, кута повороту $\theta$ або часу $t$). Тобто це одночасно і група, і гладкий геометричний простір (многовид), де перетворення можна виконувати "нескінченно малими кроками".
	
	\subsection{Алгебра Лі та комутаційні співвідношення}
	Якщо група Лі описує макроскопічні перетворення, то \textbf{алгебра Лі} описує їх поведінку поблизу одиничного елемента (в нескінченно малому околі). Елементи алгебри називаються \textbf{генераторами} (або інфінітезимальними операторами). 
	
	Ключовою властивістю алгебри Лі є те, що оператори можуть не комутувати при множенні. Алгебра Лі повністю задається \textbf{дужкою Лі} (або комутатором):
	\[ [X, Y] = X Y - Y X \]
	Комутатор показує, наскільки результат залежить від порядку застосування двох перетворень. Для будь-яких двох генераторів $X_i$ та $X_j$ їхній комутатор також належить до цієї алгебри:
	\[ [X_i, X_j] = \sum_k c_{ijk} X_k \]
	де $c_{ijk}$ — так звані \textit{структурні константи}, які однозначно визначають геометрію всієї групи.
	
	\section{Від інфінітезимального до макроскопічного}
	
	\subsection{Диференціальне рівняння еволюції}
	Нехай стан системи описується вектором $x(a)$, де $a$ — неперервний параметр перетворення. Нескінченно мала зміна стану під дією генератора $J$ має вигляд:
	\[ dx = J \, x(a) \, da \quad \implies \quad \frac{dx(a)}{da} = J x(a) \]
	Макроскопічне перетворення задається оператором (матрицею) еволюції $B(a)$, де $x(a) = B(a)x(0)$. Тоді $B(a)$ підпорядковується матричному диференціальному рівнянню:
	\[ \frac{dB(a)}{da} = J \, B(a), \quad B(0) = I \]
	
	\subsection{Мультиплікативний інтеграл та Граничний перехід}
	Еволюцію можна розглядати як послідовність нескінченної кількості нескінченно малих кроків. За методом Ейлера, розбиваючи інтервал $[0, a]$ на кроки $da$, отримуємо добуток операторів:
	\[ B(a) = \lim_{da \to 0} (I + J(a) da) \dots (I + J(da) da) (I + J(0) da) \]
	Якщо генератор $J$ є константою і не залежить від параметра $a$, цей граничний перехід дає класичну матричну експоненту:
	\[ B(a) = \exp(J a) \]
	
	Якщо ж маємо справу зі змінним генератором $J(a)$, матриці в різні моменти часу можуть не комутувати ($[J(a_1), J(a_2)] \neq 0$). Тоді застосовується \textbf{впорядкована експонента} (Path-ordered exponential), яка розписується через \textbf{ряд Дайсона} (хронологічно впорядкований інтеграл):
	\[ B(a) = I + \int_{0}^{a} J(s_1) ds_1 + \int_{0}^{a} \int_{0}^{s_1} J(s_1) J(s_2) ds_2 ds_1 + \dots \]
	
	\section{Симетрії та диференціальні рівняння}
	
	Групи Лі були створені Софусом Лі саме для вивчення диференціальних рівнянь. 
	
	Нехай маємо диференціальне рівняння, яке можна записати у вигляді дії оператора на функцію: $L(u) = 0$. 
	Група Лі є \textbf{групою симетрії} цього рівняння, якщо вона перетворює будь-який розв'язок $u$ на інший дійсний розв'язок $\tilde{u}$.
	
	У термінах алгебри Лі (генераторів) це означає наступне: генератор нескінченно малого перетворення $X$ має бути сумісним із диференціальним оператором рівняння $L$. У багатьох випадках (особливо для лінійних систем) це зводиться до умови комутації:
	\[ [X, L] = X L - L X = 0 \]
	Якщо генератор симетрії комутує з диференціальним оператором, то застосовуючи $X$ до рівняння $L(u)=0$, ми отримуємо $L(X u) = 0$. Це означає, що нова функція $(X u)$ також є розв'язком. 
	Знаючи симетрії (генератори), ми можемо знаходити нові розв'язки з уже відомих шляхом інтегрування цих генераторів. За теоремою Нетер, кожна така неперервна симетрія генерує відповідний закон збереження (наприклад, симетрія відносно зсуву в часі зберігає енергію). \textit{(Точне формулювання теореми Нетер стосується варіаційних симетрій дії (лагранжіана), а не будь-якої симетрії самого рівняння руху — але для більшості фізичних систем, що виводяться з лагранжіана, ці дві властивості збігаються.)}
	
	\section{Приклади застосування}
	
	\subsection{Квантова механіка (Рівняння Шредінгера)}
	У квантовій механіці роль генератора зсуву в часі відіграє Гамільтоніан $\hat{H}$. Рівняння еволюції хвильової функції:
	\[ \frac{\partial \psi(t)}{\partial t} = \left(-\frac{i}{\hbar}\hat{H}\right) \psi(t) \]
	Тут алгеброю Лі є антиермітові оператори. Макроскопічна еволюція системи описується оператором $U(t)$, який розраховується через впорядковану експоненту (якщо Гамільтоніан залежить від часу). Цей підхід є математичною основою для інтегралів по траєкторіях Фейнмана, де перехід між станами є сумою (інтегралом) по всіх можливих шляхах еволюції.
	
	\subsection{Обертання у 3D просторі (Група $SO(3)$)}
	Група тривимірних обертань має три генератори: $J_x, J_y, J_z$. Вони не комутують і задовольняють співвідношення алгебри Лі:
	\[ [J_x, J_y] = J_z, \quad [J_y, J_z] = J_x, \quad [J_z, J_x] = J_y \]
	(З точністю до знаків та множників). Незважаючи на наявність трьох генераторів, нескінченно мале обертання навколо довільної осі $\vec{\omega}$ задається їх лінійною комбінацією: $J_{\vec{\omega}} = \omega_x J_x + \omega_y J_y + \omega_z J_z$. Макроскопічне обертання отримується інтегруванням цього одного сумарного генератора: $R(\vec{\omega}) = \exp(J_{\vec{\omega}})$.
	
	\subsection{Приклад: Обертання 2D (Група $SO(2)$)}
	Тут існує лише один генератор $J = \left(\begin{smallmatrix} 0 & -1 \\ 1 & 0 \end{smallmatrix}\right)$.
	Пряма експонента дає класичну матрицю обертання через розклад у ряд:
	\[ \exp(J\theta) = I + J\theta + \frac{J^2\theta^2}{2!} + \dots = I \cos\theta + J \sin\theta = \begin{pmatrix} \cos\theta & -\sin\theta \\ \sin\theta & \cos\theta \end{pmatrix} \]
	Зворотне інтегрування рівняння $\frac{dB}{d\theta} = JB$ через ряд Дайсона:
	\[ B(\theta) = I + \int_{0}^{\theta} J \, ds_1 + \int_{0}^{\theta} \int_{0}^{s_1} J^2 \, ds_2 \, ds_1 + \dots \]
	де кожен наступний інтеграл від $s^{n-1}$ дає $\frac{s^n}{n}$, що точно відтворює формулу ряду Тейлора для експоненти $\sum \frac{(J\theta)^n}{n!}$.
	
	\subsection{Узагальнення на 3D: формула Родріга}
	Той самий трюк, що спрацював у прикладі $SO(2)$ вище, працює і для $SO(3)$ — лише через кубічну (а не квадратичну) періодичність генератора. Для генератора $[\omega]_\times = \omega_x J_x + \omega_y J_y + \omega_z J_z$ (лінійної комбінації з підрозділу 4.2, записаної як матриця) виконується тотожність
	\[ [\omega]_\times^3 = -\theta^2 [\omega]_\times, \qquad \theta = \|\omega\|, \]
	яка є прямим аналогом $J^2=-I$ з 2D-випадку. Підставляючи це у ряд Тейлора для експоненти, отримуємо \textbf{формулу Родріга} — явний (замкнений) вираз матриці повороту, який використовується практично в будь-якому коді для робототехніки, SLAM чи навігації:
	\[
	R = \exp\!\big([\omega]_\times\big) = I + \frac{\sin\theta}{\theta}[\omega]_\times + \frac{1-\cos\theta}{\theta^2}[\omega]_\times^2 .
	\]
	У координатах $[\omega]_\times$ — це кососиметрична матриця
	\[
	[\omega]_\times = \begin{pmatrix} 0 & -\omega_z & \omega_y \\ \omega_z & 0 & -\omega_x \\ -\omega_y & \omega_x & 0 \end{pmatrix}, \qquad [\omega]_\times v = \omega \times v \quad \forall v \in \mathbb{R}^3.
	\]
	Відображення $\omega \mapsto [\omega]_\times$ (і обернене до нього) називають операторами \textbf{«hat»} ($\wedge$) та \textbf{«vee»} ($\vee$): $\omega^\wedge = [\omega]_\times$, $\big([\omega]_\times\big)^\vee = \omega$. Це standard «переклад» між вектором кутової швидкості $\omega \in \mathbb{R}^3$ (яким оперує, наприклад, гіроскоп) та елементом алгебри $\mathfrak{so}(3)$ — і саме ця пара позначень найчастіше зустрічається в літературі з робототехніки та SLAM.
	
	\subsection{Практична альтернатива: Одиничні кватерніони}
	Хоча матриці $3\times3$ зручні для теорії, у реальному коді для дронів (наприклад, PX4, ArduPilot, ROS) орієнтацію майже завжди зберігають як \textbf{одиничний кватерніон} $q = [q_w, q_x, q_y, q_z]$. Множина одиничних кватерніонів утворює групу $S^3$, яка є подвійним покриттям $SO(3)$ (кватерніони $q$ та $-q$ задають однаковий поворот простору). Експоненційне відображення для кватерніона (яке переводить вектор $\omega \in \mathfrak{so}(3)$ у кватерніон) виглядає значно простіше та обчислюється швидше за формулу Родріга:
	\[
	\exp_q(\omega) = \begin{pmatrix} \cos(\theta/2) \\ \frac{\omega}{\theta} \sin(\theta/2) \end{pmatrix}, \quad \text{де} \quad \theta = \|\omega\|.
	\]
	Кватерніони не страждають від "шарнірного замку" (gimbal lock) і потребують менше операцій множення при композиції поворотів.
	
	\section{Практичні інструменти для робототехніки та SLAM}
	
	\subsection{Логарифмічне відображення та ретракція}
	Обернене до $\exp$ відображення — \textbf{логарифм} $\log: G \to \mathfrak{g}$ — локально відновлює генератор за заданим елементом групи (наприклад, вісь і кут повороту за матрицею $R$). Пара $(\exp, \log)$ — це «міст» між двома представленнями стану:
	\begin{itemize}
		\item алгебра $\mathfrak{g}$ (дотичний простір) — звичайний векторний простір, де коректно додавати, множити на число й усереднювати елементи;
		\item сама група $G$ — не є векторним простором: наприклад, немає коректного поняття «середнього» двох поворотів через пряме усереднення матриць, бо результат може не бути ортогональною матрицею.
	\end{itemize}
	Тому в задачах оптимізації на многовиді (SLAM, калібрування, фільтрація) стан оновлюють не додаванням, а множенням на експоненту малої поправки: $x_{\text{new}} = x \cdot \exp(\delta)$ або $x_{\text{new}} = \exp(\delta) \cdot x$, де $\delta \in \mathfrak{g}$ знаходять, наприклад, методом Гаусса--Ньютона. Такий підхід називають \textbf{ретракцією} (retraction); він автоматично гарантує, що $x_{\text{new}}$ залишається коректним елементом групи. 
	
	\textbf{Ліві та праві збурення:} У робототехніці критично важливо розрізняти, з якого боку ми множимо на $\exp(\delta)$. Праве множення $x_{\text{new}} = x \cdot \exp(\delta)$ означає поворот/зсув відносно \textit{локальної} осі робота (що природно для інтегрування даних IMU). Ліве множення $x_{\text{new}} = \exp(\delta) \cdot x$ — це перетворення відносно \textit{глобальної} (світової) системи координат.
	
	\subsection{Група $SE(3)$: обертання й зсув разом}
	Для повної пози тіла в просторі (орієнтація \emph{і} положення — саме те, що потрібно дрону чи камері) використовують групу $SE(3)$ жорстких рухів:
	\[
	T = \begin{pmatrix} R & p \\ 0 & 1 \end{pmatrix}, \quad R \in SO(3),\ p \in \mathbb{R}^3.
	\]
	Її алгебра $\mathfrak{se}(3)$ шестивимірна (три генератори обертання плюс три генератори зсуву), а елемент алгебри записують через \textit{twist}-вектор $\xi = (\omega, v) \in \mathbb{R}^6$:
	\[
	\xi^\wedge = \begin{pmatrix} [\omega]_\times & v \\ 0 & 0 \end{pmatrix} \in \mathfrak{se}(3).
	\]
	Експонента цього елемента дає не просто матрицю Родріга для обертання, а й «перекручену» версію зсуву:
	\[
	\exp(\xi^\wedge) = \begin{pmatrix} R & Vv \\ 0 & 1 \end{pmatrix}, \qquad V = I + \frac{1-\cos\theta}{\theta^2}[\omega]_\times + \frac{\theta-\sin\theta}{\theta^3}[\omega]_\times^2,
	\]
	де $R$ — та сама матриця з підрозділу 4.4, а $V$ (\textit{лівий якобіан} $SO(3)$) враховує, що зсув відбувається \emph{одночасно} з обертанням, а не після нього.
	
	\subsection{Приєднане представлення (Adjoint)}
	\textbf{Приєднане представлення} $\mathrm{Ad}_T : \mathfrak{g} \to \mathfrak{g}$ показує, як елемент алгебри (наприклад, швидкість або збурення) перетворюється при зміні системи відліку (матричний запис: $\mathrm{Ad}_T(X) = T X T^{-1}$). 
	
	Для $SO(3)$ в термінах векторів кутової швидкості це зводиться до $\mathrm{Ad}_R(\omega) = R\omega$ — обертання самого вектора. Практичне значення: якщо є похибка (коваріація) оцінки орієнтації в одній системі координат (наприклад, у тілі дрона), $\mathrm{Ad}$ дозволяє коректно перенести цю похибку в іншу систему координат (світову) — саме це відбувається «під капотом» у фільтрах Калмана на многовидах (Invariant EKF) і в графових SLAM-оптимізаторах.
	
	\subsection{Формула Бейкера--Кемпбелла--Хаусдорфа (BCH)}
	Оскільки генератори зазвичай не комутують, добуток двох експонент не дорівнює експоненті суми: $\exp(X)\exp(Y) \neq \exp(X+Y)$. Точний зв'язок дає формула BCH:
	\[
	\log\big(\exp(X)\exp(Y)\big) = X + Y + \tfrac{1}{2}[X,Y] + \tfrac{1}{12}\big([X,[X,Y]] + [Y,[Y,X]]\big) + \dots
	\]
	Перший поправковий доданок $[X,Y]$ і є причина, чому \emph{порядок} застосування двох малих поворотів має значення. На практиці (з наближенням першого порядку) цю формулу використовують, щоб об'єднувати послідовні малі приростки пози в один сумарний, не перераховуючи все спочатку.
	
	\section{Де це застосовується: розпізнавання, SLAM, навігація дронів}
	
	\subsection{Еквіваріантні нейронні мережі (розпізнавання)}
	Звичайна згорткова мережа (CNN) за побудовою інваріантна лише до зсувів (translation) — поворот об'єкта на вхідному зображенні змушує мережу вчити цю симетрію «з нуля» на даних. \textbf{Group-equivariant CNN} (Cohen, Welling, ICML 2016) вбудовують ширшу групу симетрій (повороти, віддзеркалення) безпосередньо в архітектуру згортки, тож мережа гарантовано узгоджено реагує на повернутий вхід без додаткової аугментації даних. \textbf{LieConv} (Finzi, Stanton, Izmailov, Wilson, ICML 2020) узагальнює цю ідею на довільну групу Лі, для якої достатньо задати лише $\exp$ і $\log$ — тобто саме той апарат, що описаний у підрозділі 5.1 — і працює навіть на нерегулярних даних (наприклад, хмарах точок з лідара чи радара дрона).
	
	\subsection{SLAM: оцінка пози на многовиді}
	У сучасних SLAM-системах (граф поз, bundle adjustment — наприклад, у бібліотеках на кшталт GTSAM чи g2o) орієнтація й поза камери або робота зберігаються не як довільна трійка кутів, а як елемент $SO(3)$ або $SE(3)$. Оптимізація (пошук пози, що найкраще узгоджує всі виміри) виконується через ретракцію з підрозділу 5.1: на кожній ітерації знаходять малу поправку $\delta$ в алгебрі й оновлюють позу множенням на $\exp(\delta)$, а не додаванням. Приєднане представлення (5.3) при цьому переноситиме коваріацію похибки між локальною системою відліку сенсора та глобальною картою.
	
	\subsection{Навігація дронів: інтегрування орієнтації гіроскопа}
	Це пряме продовження рівняння $\dot x = Jx$ із розділу 2: кутова швидкість $\omega(t)$, яку вимірює гіроскоп, і є генератор $J(t)$, а орієнтація дрона $R(t)$ — це те саме макроскопічне перетворення $B(a)$. Оскільки $\omega(t)$ змінюється з часом, а виміри надходять на високій частоті (сотні--тисячі разів на секунду), пряме інтегрування рівняння еволюції на кожному кроці занадто дороге. \textbf{IMU preintegration} (Forster, Carlone, Dellaert, Scaramuzza, RSS 2015 / IEEE T-RO 2017) розв'язує це, використовуючи формулу BCH (5.4), щоб заздалегідь «згорнути» усі виміри гіроскопа й акселерометра між двома ключовими кадрами в один сумарний приріст пози на $SE(3)$ — це один із базових прийомів практично в кожній сучасній системі візуально-інерціальної навігації дрона.
	
\end{document}